\documentclass[12pt]{article}

\usepackage[margin=2.5cm]{geometry}
\usepackage{tikz}
\usepackage{amsmath}
\usepackage{mathtools}
\usepackage{graphicx}

\setlength{\parindent}{0pt}
\setlength{\parskip}{0.75em}


\title{Chapter 8: Gravitation}
\author{İskender Yalçınkaya, PhD}
\date{\today}

\begin{document}
\maketitle




\section{Worked Examples}

\subsection{Newton's Discovery of the Inverse-Square Law}

In the 17th century, Johannes Kepler established that the orbital
period $T$ of a planet and the semi-major axis $a$ of its elliptical
orbit satisfy Kepler's third law,
\begin{equation}
    T^2 \propto a^3.
\end{equation}

According to his later account, Isaac Newton used Kepler's third
law and the dynamics of circular motion to infer how the force
maintaining a planet's orbit depends on its distance from the Sun.

We will reconstruct Newton's reasoning by approximating planetary
orbits as circles. Consider a planet of mass $m$ moving in a circular
orbit of radius $r$ around the Sun, completing one revolution in
a period $T$.

\begin{enumerate}
    \item[(a)] Express the orbital speed $v$ of the planet in terms
    of $r$ and $T$.

    \item[(b)] Using Newton's second law, determine the magnitude
    of the force required to maintain the planet's circular motion.

    \item[(c)] Using Kepler's third law, show that the force
    maintaining the planet's orbit is inversely proportional
    to the square of its distance from the Sun,
    \begin{equation}
        F(r)\propto\frac{1}{r^2},
    \end{equation}
    for a planet of fixed mass.

    \item[(d)] Newton also investigated whether the force that
    causes objects to fall toward Earth could be the same force
    that keeps the Moon in its orbit.

    Assume that Earth's gravitational attraction follows the same
    inverse-square law. The Moon is approximately $60$ Earth radii
    from Earth's center. Determine the gravitational acceleration
    of the Moon in terms of the gravitational acceleration $g$
    at Earth's surface.

    Compare your result with the centripetal acceleration of
    the Moon, assuming a circular orbit with a period of
    $T_M=27.3$ days.

    Use $R_E=6.37\times10^6\,\mathrm{m}$ and
    $g=9.81\,\mathrm{m/s^2}$.
\end{enumerate}

\textbf{Solution:}

\begin{enumerate}
    \item[(a)] Since the planet completes one full revolution
    in a period $T$, it travels a distance equal to the
    circumference of its circular orbit, $2\pi r$. Therefore,
    its orbital speed is
    \begin{equation}
        \boxed{v=\frac{2\pi r}{T}.}
    \end{equation}

    \item[(b)] The planet undergoes uniform circular motion
    with centripetal acceleration
    \begin{equation}
        a_c=\frac{v^2}{r}.
    \end{equation}
    By Newton's second law, the magnitude of the force
    responsible for this acceleration is
    \begin{equation}
        F=ma_c=\frac{mv^2}{r}.
    \end{equation}
    Substituting the expression for the orbital speed gives
    \begin{equation}
        F=\frac{m}{r}\left(\frac{2\pi r}{T}\right)^2.
    \end{equation}
    Simplifying,
    \begin{equation}
        \boxed{F=\frac{4\pi^2mr}{T^2}.}
    \end{equation}

    \item[(c)] For circular orbits, the semi-major axis is
    equal to the orbital radius, $a=r$. Therefore, Kepler's
    third law can be written as
    \begin{equation}
        T^2=Cr^3,
    \end{equation}
    where $C$ is a constant that is the same for all planets
    orbiting the Sun.

    Substituting this expression into the force obtained in
    part (b), we find
    \begin{equation}
        F=\frac{4\pi^2mr}{Cr^3}
        =\frac{4\pi^2m}{Cr^2}.
    \end{equation}
    Since $4\pi^2/C$ is constant, we conclude that
    \begin{equation}
        \boxed{F(r)\propto\frac{1}{r^2}}
    \end{equation}
    for a planet of fixed mass.

    Thus, the force required to maintain planetary orbits
    decreases with the square of the distance from the Sun.

    \item[(d)] Let $M_E$ denote Earth's mass. Assuming that
    Earth's gravitational attraction obeys the inverse-square
    law, the gravitational acceleration at a distance $r$
    from Earth's center has the form
    \begin{equation}
        g(r)=\frac{K}{r^2},
    \end{equation}
    where $K$ is a constant determined by Earth.

    At Earth's surface, $r=R_E$, so
    \begin{equation}
        g=\frac{K}{R_E^2}.
    \end{equation}
    At the Moon's orbital radius, $r_M=60R_E$,
    the gravitational acceleration is therefore
    \begin{equation}
        g_M=\frac{K}{(60R_E)^2}
        =\frac{1}{60^2}\frac{K}{R_E^2}.
    \end{equation}
    Hence,
    \begin{equation}
        \boxed{g_M=\frac{g}{3600}.}
    \end{equation}

    Using $g=9.81\,\mathrm{m/s^2}$, we obtain
    \begin{equation}
        \boxed{g_M\approx2.73\times10^{-3}\,\mathrm{m/s^2}.}
    \end{equation}

    We can now compare this prediction with the centripetal
    acceleration required to maintain the Moon's orbit.

    The Moon's orbital radius is approximately
    \begin{equation}
        r_M=60R_E=3.822\times10^8\,\mathrm{m},
    \end{equation}
    and its orbital period in seconds is
    \begin{equation}
        T_M=27.3\times24\times3600
        =2.35872\times10^6\,\mathrm{s}.
    \end{equation}

    The centripetal acceleration is
    \begin{equation}
        a_M=\frac{v_M^2}{r_M}
        =\frac{4\pi^2r_M}{T_M^2}.
    \end{equation}
    Substituting the numerical values gives
    \begin{equation}
        \boxed{a_M\approx2.71\times10^{-3}\,\mathrm{m/s^2}.}
    \end{equation}

    We therefore find
    \begin{equation}
        \boxed{a_M\approx g_M.}
    \end{equation}

    The agreement between these two independently calculated
    accelerations supports the hypothesis that the force
    responsible for the Moon's orbital motion is the same
    gravitational attraction that causes objects to fall
    toward Earth's surface.
\end{enumerate}

\end{document}
